\section{(H) 函数}

命题 3. 若 \(\mathfrak{K}_0\) 是一个哈代域，则存在一个哈代域 \(\mathfrak{K}\) （包含 \(\mathfrak{K}_0\) ），使得对每个函数 \(z \in \mathfrak{K}\) （在 \(+\infty\) 的某个邻域上不恒为零）， \(e^z\) 与 \(\log |z|\) 都属于 \(\mathfrak{K}\) 。

用 \(\mathfrak{K}\) 表示具有下列性质的函数 \(f\in \mathcal{H}(\mathfrak{F},\mathbf{R})\) 所成的集：对每个函数 \(f\in \mathfrak{K}\) ，存在有限个哈代域 \(\mathfrak{K}_1,\mathfrak{K}_2,\ldots ,\mathfrak{K}_n\) （数 \(n\) 与诸域 \(\mathfrak{K}_i\) 依赖于 \(f\) ），使得 \(f\in \mathfrak{K}_n\) ，且对 \(0\leqslant i\leqslant n - 1\) 有 \(\mathfrak{K}_{i + 1} = \mathfrak{K}_i(u_{i + 1})\) ，其中 \(u_{i + 1}\) 等于 \(e^{z_i}\) 或 \(\log |z_i|\) ，而 \(z_{i}\) 属于 \(\mathfrak{K}_i\) 且在 \(+\infty\) 的某个邻域上不恒为零。称 \(u_{1},u_{2},\ldots ,u_{n}\) 构成域 \(\mathfrak{K}_n\) 以及函数 \(f\) 的一个定义序列；同一个函数自然可以有多个定义序列。

由 V, p.~247 的定义 1，每个函数 \(f \in \mathfrak{K}\) ，若在 \(+\infty\) 的某个邻域上不恒为零，则在某个区间 \([x_0, +\infty[\) 上具有常号且可导；若 \(f \in \mathfrak{K}_n\) ，则 \(1/f\) 与 \(f'\) 等于 \(\mathfrak{K}_n\) 中的函数、从而 \(\mathfrak{K}\) 中的函数（在 \(+\infty\) 的某个邻域上）。要看出 \(\mathfrak{K}\) 是一个哈代域，只需证明：若 \(f\) 与 \(g\) 是 \(\mathfrak{K}\) 中两个函数，则 \(f - g\) 与 \(fg\) 等于 \(\mathfrak{K}\) 中的函数（在 \(+\infty\) 的某个邻域上）。设 \(u_1, u_2, \ldots, u_m\) 是 \(f\) 的一个定义序列， \(v_1, v_2, \ldots, v_n\) 是 \(g\) 的一个定义序列。序列 \(u_1, u_2, \ldots, u_m, v_1, v_2, \ldots, v_n\) （由 \((u_i)\) 与 \((v_i)\) 拼接而得）仍是某个哈代域 \(\mathfrak{K}_{m+n}\) 的定义序列，而这个域包含 \(f\) 与 \(g\) ，故 \(f - g\) 与 \(fg\) 等于 \(\mathfrak{K}_{m+n}\) 中的函数（在 \(+\infty\) 的某个邻域上）。

称在命题 3 的证明中定义的哈代域 \(\mathfrak{K}\) 为哈代域 \(\mathfrak{K}_0\) 的 (H) 扩张。

若 \(\mathfrak{K}'\) 是另一个具有命题 3 所述性质的哈代域，则由 \(\mathfrak{K}\) 的构造可知 \(\mathfrak{K} / \mathrm{R}_{\infty}\) 包含于 \(\mathfrak{K}' / \mathrm{R}_{\infty}\) 。（滥用语言的说法）称哈代域 \(\mathfrak{K}_0\) 的 (H) 扩张是具有这些性质的最小哈代域 \(\mathfrak{K}\) 。

定义 2. 称实系数有理函数的哈代域 \(\mathbf{R}(x)\) 的 (H) 扩张为 (H) 函数域。属于这个扩张的任一函数称为一个 (H) 函数。

按照这个定义，若 f 是一个 (H) 函数，在 \(+\infty\) 的某个邻域上不恒为零，则 \(e^{f}\) 与 \(\log |f|\) 也是 (H) 函数。更一般地，若 g 是第二个 (H) 函数， \(u_{1}, u_{2}, \ldots, u_{n}\) 是 g 的一个定义序列，且 \(f(x)\) 随 x 趋于 \(+\infty\) ，则对 n 作归纳可见复合函数 \(u_{1} \circ f, u_{2} \circ f, \ldots, u_{n} \circ f\) 与 \(g \circ f\) 都是 (H) 函数。

